\documentclass[11pt,letterpaper,twocolumn]{article}

% ── ACL-style formatting ─────────────────────────────────────────────────────
\usepackage{times}
\usepackage{microtype}
\usepackage[margin=1in,columnsep=0.25in]{geometry}
\usepackage{titlesec}
\usepackage{booktabs}
\usepackage{amsmath,amssymb}
\usepackage{graphicx}
\usepackage{enumitem}
\usepackage{url}
\usepackage[hidelinks]{hyperref}
\usepackage{natbib}
\usepackage{filecontents}
\usepackage{authblk}
\usepackage{caption}
\usepackage{array}
\usepackage{multirow}

% ACL-style section headers
\titleformat{\section}{\normalfont\bfseries}{\thesection}{0.5em}{}
\titleformat{\subsection}{\normalfont\bfseries\itshape}{\thesubsection}{0.5em}{}
\titlespacing{\section}{0pt}{6pt}{2pt}
\titlespacing{\subsection}{0pt}{4pt}{1pt}

\setlength{\parindent}{0pt}
\setlength{\parskip}{4pt}
\captionsetup{font=small}

% ── Bib ───────────────────────────────────────────────────────────────────────
\begin{filecontents}{refs_progress.bib}
@inproceedings{chimani2022spanner,
  author    = {Markus Chimani and Finn Stutzenstein},
  title     = {Spanner Approximations in Practice},
  booktitle = {30th Annual European Symposium on Algorithms ({ESA} 2022)},
  series    = {{LIPIcs}},
  volume    = {244},
  pages     = {37:1--37:15},
  year      = {2022},
  doi       = {10.4230/LIPIcs.ESA.2022.37}
}
\end{filecontents}

% ── Title ─────────────────────────────────────────────────────────────────────
\title{\textbf{Graph Spanner Approximations in Practice:\\
Interim Progress Report}}

\author[1]{Sheikh Muhammad Muneeb \quad Abdul Moiz Ahmed \quad Moosa Owais}
\affil[1]{Habib University, Karachi \quad
\small\texttt{\{ms08914, aa08811, mo08952\}@st.habib.edu.pk}}
\date{April 2026}

\begin{document}
\maketitle

% ── Abstract ─────────────────────────────────────────────────────────────────
\begin{abstract}
We report interim progress on implementing and evaluating the Greedy
$(2k{-}1)$-spanner algorithm from \citet{chimani2022spanner}.
Our implementation produces spanners that satisfy the stretch
guarantee on all tested instances.
Across 5 graph sizes ($n \in \{10, 20, 50, 100, 200\}$),
3 density levels, and 2 stretch parameters ($\alpha \in \{3, 5\}$),
we observe that sparseness scales sub-linearly with $n$,
effective stretch consistently stays below the theoretical bound,
and runtime grows steeply with density for $\alpha=3$.
Baswana--Sen implementation and full-scale benchmarking remain for
the final week.
\end{abstract}

% ── 1. Problem ───────────────────────────────────────────────────────────────
\section{Problem Statement}

A \emph{multiplicative $\alpha$-spanner} $H$ of a weighted undirected
graph $G=(V,E,w)$ is a sparse subgraph satisfying
$d_H(u,v) \leq \alpha \cdot d_G(u,v)$ for all $u,v \in V$.
The fundamental tension is between \emph{stretch} $\alpha$ and
\emph{sparseness} $|E_H|/n$: a smaller $\alpha$ forces more edges into $H$,
increasing both size and construction time.

This problem is NP-hard to solve optimally.
\citet{chimani2022spanner} provide the first systematic experimental
comparison of polynomial-time approximation algorithms, evaluating five
algorithms across 2,000+ graph instances.
Our project implements and benchmarks the two most practical algorithms
from that study: the \textbf{Greedy} algorithm (ADDJS) and the
\textbf{Baswana--Sen} (BS) clustering algorithm.

% ── 2. Algorithm ─────────────────────────────────────────────────────────────
\section{Algorithm: Greedy Spanner}

The Greedy $(2k{-}1)$-spanner processes edges of $G$ in non-decreasing
weight order.
Each edge $(u,v)$ is added to the spanner $H$ only if no path
already exists in $H$ satisfying the stretch bound:

\begin{equation}
d_H(u,v) \;>\; (2k-1)\cdot w(u,v)
\end{equation}

Each membership test runs Dijkstra on the current (sparse, growing) $H$.
The algorithm is guaranteed to produce a spanner of size
$O(n^{1+1/k})$ with lightness $O(n/k)$.

\paragraph{Implementation.}
Our Python implementation uses NetworkX for graph primitives and
Dijkstra shortest-path queries.
Correctness is verified by checking $d_H(u,v) \leq \alpha \cdot d_G(u,v)$
for all pairs (exact for $n \leq 150$, 500-sample estimate for larger).

% ── 3. Preliminary Results ───────────────────────────────────────────────────
\section{Preliminary Results}

We ran the Greedy algorithm on 90 Erd\H{o}s--R\'enyi instances:
$n \in \{10,20,50,100,200\}$, density $\in \{0.1, 0.3, 0.5\}$,
3 graphs per combination, $k \in \{2,3\}$.
Table~\ref{tab:results} reports means across the 3 repetitions.
All effective-stretch values are at or below the theoretical bound
$\alpha = 2k-1$, confirming correctness.

\begin{table}[t]
\centering
\small
\setlength{\tabcolsep}{3pt}
\begin{tabular}{r r r r r r r}
\toprule
$n$ & dens & $\alpha$ & sparse & light & stretch & time(s) \\
\midrule
 10 & 0.1 & 3 & 0.90 & 1.00 & 1.03 & 0.001 \\
 10 & 0.3 & 3 & 0.90 & 1.00 & 2.20 & 0.001 \\
 50 & 0.1 & 3 & 1.13 & 1.25 & 2.88 & 0.008 \\
 50 & 0.3 & 3 & 1.16 & 1.31 & \textbf{3.00} & 0.036 \\
 50 & 0.5 & 3 & 1.16 & 1.32 & \textbf{3.00} & 0.059 \\
100 & 0.1 & 3 & 1.18 & 1.35 & 2.99 & 0.071 \\
100 & 0.3 & 3 & 1.20 & 1.37 & \textbf{3.00} & 0.271 \\
100 & 0.5 & 3 & 1.29 & 1.52 & \textbf{3.00} & 0.491 \\
200 & 0.1 & 3 & 1.27 & 1.52 & 2.72 & 0.719 \\
200 & 0.3 & 3 & 1.34 & 1.65 & 2.76 & 2.243 \\
200 & 0.5 & 3 & 1.34 & 1.60 & 2.62 & 3.822 \\
\midrule
 50 & 0.3 & 5 & 1.01 & 1.05 & 4.28 & 0.026 \\
100 & 0.3 & 5 & 1.03 & 1.05 & 4.76 & 0.243 \\
200 & 0.3 & 5 & 1.09 & 1.15 & 4.37 & 2.147 \\
\bottomrule
\end{tabular}
\caption{Greedy results (mean over 3 runs). \textbf{Bold}: stretch tight at bound.}
\label{tab:results}
\end{table}

\paragraph{Key observations.}

\textbf{(1) Sparseness grows sub-linearly.}
For $\alpha=3$, sparseness rises from 0.90 at $n=10$ to 1.34 at
$n=200$, consistent with the $O(n^{1+1/k}) = O(n^{3/2})$ edge-count
guarantee (in terms of edges-per-vertex the growth is slow).

\textbf{(2) Higher stretch $\Rightarrow$ sparser spanners.}
At $n=200$, density 0.3: $\alpha=3$ gives sparseness 1.34 vs.\
$\alpha=5$ gives 1.09, a 19\% reduction in edges per vertex.

\textbf{(3) Stretch bound is tight at medium-to-high density.}
Several instances at density 0.3 and 0.5 achieve effective stretch
exactly equal to $\alpha$, confirming our correctness oracle.

\textbf{(4) Runtime grows sharply with density.}
At $n=200$: density 0.1 takes 0.72\,s; density 0.5 takes 3.82\,s ---
a $5.3\times$ increase driven by the quadratically larger edge set
requiring Dijkstra per edge.

\begin{figure}[t]
\centering
\includegraphics[width=\columnwidth]{/mnt/user-data/outputs/spanner_results/fig1_sparseness.png}
\caption{Sparseness vs.\ $n$ for $\alpha=3$ (left) and $\alpha=5$ (right).
Higher density yields higher sparseness; $\alpha=5$ is consistently sparser.}
\label{fig:sparse}
\end{figure}

% ── 4. Challenges ────────────────────────────────────────────────────────────
\section{Challenges Identified}

\paragraph{Greedy scalability.}
The Dijkstra-per-edge design is prohibitive for large dense graphs.
At $n=200$, density 0.5, runtime is already 3.8\,s.
Extrapolating to $n=2{,}000$ (the paper's largest Random instance)
at moderate density would take hours.
Mitigation: restrict Greedy benchmarks to $n \leq 500$ and
sparse graphs ($\text{density} \leq 0.3$), matching the original
paper's reported solve-rate limitations.

\paragraph{Stretch oracle at scale.}
All-pairs shortest-path verification is $O(n^2 \log n)$ and
infeasible beyond $n \approx 500$.
Our 500-pair sampling approximation introduces a small probability
of underestimating the true maximum stretch.

\paragraph{Baswana--Sen correctness.}
BS involves multi-round cluster sampling with randomized cluster
membership.
The key correctness invariant (every vertex is either clustered
or has added its inter-cluster edges) requires careful bookkeeping
distinct from the Greedy case.

% ── 5. Plan ──────────────────────────────────────────────────────────────────
\section{Next Steps and Timeline}

\paragraph{Week 13 (this week).}
Implement Baswana--Sen algorithm.
Run preliminary BS vs.\ Greedy comparison on Random instances
($n \leq 500$) for $k \in \{2,3\}$.

\paragraph{Week 14.}
Extend experiments to TSPLIB (95 instances, $n$ up to 1,200 for
Greedy; all for BS) and the SNAP Facebook/Gnutella graphs.
Measure and compare: sparseness, lightness, effective stretch
(sampled), peak memory, and runtime across all five metrics.

\paragraph{Week 15.}
Analyse results, produce final figures and tables replicating
the structure of \citet{chimani2022spanner}, and prepare the
final presentation.

% ── References ───────────────────────────────────────────────────────────────
\bibliographystyle{plainnat}
\bibliography{refs_progress}

\end{document}
